\subsection{\texorpdfstring{INNER PRODUCTS BETWEEN \(\mathsf{T}(\mathbf{M})\) AND \(\mathsf{T}(\mathbf{M}^{*})\), \(\mathsf{S}(\mathbf{M})\) AND \(\mathsf{S}(\mathbf{M}^{*})\), \(\bigwedge(\mathbf{M})\) AND \(\bigwedge(\mathbf{M}^{*})\)}{Inner products between T(M) and T(M*), S(M) and S(M*), wedge(M) and wedge(M*)}}

The right inner product defines on \(\mathsf{T}(\mathbf{M})\) (resp. \(\mathsf{S}(\mathbf{M})\) , resp. \(\bigwedge (\mathbf{M})\) ) a right module structure over the algebra \(\mathsf{T}(\mathbf{M})^{*\mathfrak{g}\mathfrak{r}}\) (resp. \(\mathsf{S}(\mathbf{M})^{*\mathfrak{g}\mathfrak{r}}\) , resp. \(\bigwedge (\mathbf{M})^{*\mathfrak{g}\mathfrak{r}}\) ) (no. 7, Examples). Using the canonical homomorphisms \(\theta_{\mathsf{T}}\) (resp. \(\theta_{\mathsf{S}}\) , resp. \(\theta_{\wedge}\) ) of no. 5, we derive

a right \(\mathsf{T}(\mathbf{M}^{*})\) -module structure on \(\mathsf{T}(\mathbf{M})\)

a right \(\mathbf{S}(\mathbf{M}^{*})\) -module structure on \(\mathbf{S}(\mathbf{M})\)

a left \(\bigwedge (\mathbf{M}^{*})\) -module structure on \(\bigwedge (\mathbf{M})\) .

The external law of any of these structures is also denoted by

\[
(z ^ {*}, t) \mapsto i (z ^ {*}). t
\]

(by an abuse of language); we also write \(t \llcorner z^*\) instead of \(i(z^*)\) . \(t\) in the case of \(\mathsf{T}(\mathbf{M})\) or \(\mathsf{S}(\mathbf{M})\) ; on the other hand, we write \(z^* \lrcorner t\) in the case of \(\bigwedge (\mathbf{M})\) and say that this is a left inner product of \(t\) by \(z^*\) , since then we have a left \(\bigwedge (\mathbf{M}^*)\) -module law. For \(z^*\) homogeneous of degree \(n\) and \(t\) homogeneous of degree \(p\) , \(i(z^*)\) . \(t = 0\) if \(p < n\) and, for \(x_i \in \mathbf{M} (1 \leqslant i \leqslant p)\) , \(x_j^* \in \mathbf{M}^* (1 \leqslant j \leqslant n)\) and \(p \geqslant n\) , by virtue of formulae (58), (59) and (60) of no. 7,

\[
\begin{array}{r l} i (x _ {1} ^ {*} \otimes x _ {2} ^ {*} \otimes \dots \otimes x _ {n} ^ {*}) \cdot (x _ {1} \otimes x _ {2} \otimes \dots \otimes x _ {p}) \\ & = \left(\prod_ {j = 1} ^ {n} \langle x _ {j} ^ {*}, x _ {j} \rangle\right) x _ {n + 1} \otimes \dots \otimes x _ {p} \end{array}\tag{66}
\]

\[
i \left(x _ {1} ^ {*} x _ {2} ^ {*} \dots x _ {n} ^ {*}\right). (x _ {1} x _ {2} \dots x _ {p}) = \sum_ {\sigma} \left(\prod_ {j = 1} ^ {n} \langle x _ {j} ^ {*}, x _ {\sigma (j)} \rangle\right) x _ {\sigma (n + 1)} \dots x _ {\sigma (p)}\tag{67}
\]

\[
\begin{array}{l} i \left(x _ {1} ^ {*} \wedge x _ {2} ^ {*} \wedge \dots \wedge x _ {n} ^ {*}\right). \left(x _ {1} \wedge x _ {2} \wedge \dots \wedge x _ {p}\right) \\ = (- 1) ^ {n (n - 1) / 2} \sum_ {\sigma} \varepsilon_ {\sigma} \left(\prod_ {j = 1} ^ {n} \langle x _ {j} ^ {*}, x _ {\sigma (j)} \rangle\right) x _ {\sigma (n + 1)} \wedge \dots \wedge x _ {\sigma (p)} \end{array}\tag{68}
\]

where, the formulae (67) and (68), \(\sigma\) runs through the set of permutations \(\sigma \in \mathfrak{S}_p\) which are increasing on the intervals \([1, n]\) and \([n + 1, p]\) . We can also write, in the inner product notation,

\[
\begin{array}{l l} \langle t \llcorner u ^ {*}, v ^ {*} \rangle = \langle t, \theta_ {\mathsf {T}} (u ^ {*} v ^ {*}) \rangle & \text {for} \quad t \in \mathsf {T} (\mathbf {M}), u ^ {*}, v ^ {*} \quad \text {in} \quad \mathsf {T} (\mathbf {M} ^ {*}) \\ \langle t \llcorner u ^ {*}, v ^ {*} \rangle = \langle t, \theta_ {\mathsf {S}} (u ^ {*} v ^ {*}) \rangle & \text {for} \quad t \in \mathsf {S} (\mathbf {M}), u ^ {*}, v ^ {*} \quad \text {in} \quad \mathsf {S} (\mathbf {M} ^ {*}) \\ \langle v ^ {*}, u ^ {*} \lrcorner t \rangle = \langle \theta_ {\wedge} (u ^ {*} \wedge v ^ {*}), t \rangle & \text {for} \quad t \in \bigwedge (\mathbf {M}), u ^ {*}, v ^ {*} \quad \text {in} \quad \bigwedge (\mathbf {M} ^ {*}). \end{array}
\]

We leave to the reader the task of finding explicitly the analogous formulae for left inner products, this time using formulae (61), (62) and (63).

The above can be applied with M replaced by its dual \(M^{*}\) ; \(M^{*}\) must then be replaced by the bidual \(M^{**}\) and \(T(M^{*})\) , for example, thus has a right module structure over the algebra \(T(M^{**})\) . But the canonical mapping \(c_{M}: M \to M^{**}\) defines an algebra homomorphism \(T(c_{M}): T(M) \to T(M^{**})\) , by means of which \(T(M^{*})\) has a right \(T(M)\) -module structure. Similarly \(S(M^{*})\) (resp. \(\bigwedge(M^{*})\) ) has a right \(S(M)\) -module (resp. left \(\bigwedge(M)\) -module) structure. The explicit formulae giving the external laws of these modules are derived immediately from the above by exchanging the roles of M and \(M^{*}\) . Note that, for all \(x \in M\) , \(i(x)\) is always a derivation (resp. antiderivation of zero square) of the graded algebra \(S(M^{*})\) (resp. \(\bigwedge(M^{*})\) ).

PROPOSITION 11. The canonical homomorphism \(\theta_{\mathsf{T}}: \mathsf{T}(\mathbf{M}^{*}) \to \mathsf{T}(\mathbf{M})^{*\mathsf{gr}}\) (resp. \(\theta_{\mathsf{S}}: \mathsf{S}(\mathbf{M}) \to \mathsf{S}(\mathbf{M})^{*\mathsf{gr}}\) , resp. \(\theta_{\wedge}: \bigwedge (\mathbf{M}^{*}) \to \bigwedge (\mathbf{M})^{*\mathsf{gr}}\) ) is a right \(\mathsf{T}(\mathbf{M})\) -module (resp. right \(\mathsf{S}(\mathbf{M})\) -module, resp. left \(\bigwedge (\mathbf{M})\) -module) homomorphism.

We show first that, for \(z^{*} \in \mathsf{T}(\mathbf{M}^{*})\) and \(t \in \mathsf{T}(\mathbf{M})\) ,

\[
\theta_ {T} (z ^ {*} \llcorner t) = \theta_ {T} (z ^ {*}) \llcorner t.\tag{69}
\]

Since M is a generating system of the algebra \(\mathsf{T}(\mathbf{M})\) , we need only prove (69) when \(t = x \in M\) ; moreover we can restrict our attention to the case where \(z^{*} = x_{1}^{*} \otimes x_{2}^{*} \otimes \cdots \otimes x_{p}^{*}\) , where the \(x_{j}^{*} \in M^{*}\) , and then, by (66) with the roles of M and \(M^{*}\) interchanged, \(z^{*} \llcorner x = \langle x, x_{1}^{*} \rangle x_{2}^{*} \otimes \cdots \otimes x_{p}^{*}\) . Therefore, for all \(y_{2}, \ldots, y_{p}\) in M,

\[
\begin{array}{r l} \langle \theta_ {T} (z ^ {*} \llcorner x), y _ {2} \otimes y _ {3} \otimes \dots \otimes y _ {p} \rangle & = \langle x, x _ {1} ^ {*} \rangle \prod_ {j = 2} ^ {p} \langle y _ {j}, x _ {j} ^ {*} \rangle \\ & = \langle \theta_ {T} (z ^ {*}), x \otimes y _ {2} \otimes \dots \otimes y _ {p} \rangle \\ & = \langle \theta_ {T} (z ^ {*}) \llcorner x, y _ {2} \otimes \dots \otimes y _ {p} \rangle \end{array}
\]

whence (69).

We prove secondly that for \(z^{*} \in \mathbf{S}(\mathbf{M}^{*})\) and \(t \in \mathbf{S}(\mathbf{M})\) ,

\[
\theta_ {S} (z ^ {*} \llcorner t) = \theta_ {S} (z ^ {*}) \llcorner t.\tag{70}
\]

As above, we can limit ourselves to the case \(t = x \in \mathbf{M}\) . But further, here \(i(x)\) is a derivation of \(\mathbf{S}(\mathbf{M}^*)\) and a derivation of \(\mathbf{S}(\mathbf{M})^{*\text{gr}}\) . Therefore (§ 10, no. 7, Corollary to Proposition 9) it suffices to verify (70) for \(z^* = x^* \in \mathbf{M}^*\) , since \(\mathbf{M}^*\) is a generating system of \(\mathbf{S}(\mathbf{M}^*)\) ; but this is trivial, the two sides then being equal to \(\langle x^*, x \rangle\) . A similar argument proves the relation

\[
\theta_ {\wedge} (t \lrcorner z ^ {*}) = t \lrcorner \theta_ {\wedge} (z ^ {*})\tag{71}
\]

for \(z^{*} \in \bigwedge(\mathbf{M}^{*})\) and \(t \in \bigwedge(\mathbf{M})\) : observe then that, for \(x \in \mathbf{M}\) , \(i(x)\) is an antiderivation in \(\bigwedge(\mathbf{M}^{*})\) as well as in \(\bigwedge(\mathbf{M})^{*\text{er}}\) and use § 10, no. 7, Corollary to Proposition 9. There is an analogous result for left inner products.

